% BEGIN THERMAL_R03_SYNC_01
\section{Introducción}

\subsection{Determinación de las constantes y relaciones entre acción, área e información}

Este artículo determina las secciones de acción de Planck y su retorno,
el funcional de Barbero--Immirzi y las coordenadas angulares CKM, y
establece la transducción de Boltzmann y la reciprocidad gravitatoria en
las realizaciones dimensionales especificadas. Son resultados vinculados
por operaciones efectivas: la acción convierte frecuencia en energía;
Barbero--Immirzi pondera la incidencia que resuelve el área; Boltzmann
convierte información adimensional en entropía; y la gravitación
transforma la acción en una escala geométrica común. La cuestión física
es cómo estas funciones y sus valores se determinan desde una misma
estructura discreta, en lugar de entrar como parámetros medidos.

La determinación de Planck separa el orden decimal de la acción de sus
coordenadas de retorno. El cociente local del registro posee dieciséis
hojas; su norma multiplicativa, seguida de la autoescala, da
\(\operatorname{ord}_{10}(\varphi\alpha^{16})=-34\).
Los caracteres regionales fijan después
\(\hbar_{\rm pre}=H_5\,10^{-34}\mathcal U_S\) y
\(\hbar_{\rm ret}=\eta_{\rm ret}\,10^{-34}\mathcal U_S\),
con \(0<\eta_{\rm ret}<H_5\). El cociente
\(R_{\rm act}=H_5/\eta_{\rm ret}>1\) conserva el efecto del
retorno al cambiar la base positiva común; en cada sección,
\(h_a=2\pi\hbar_a\). El lector local, el carácter y la regla de
continuación regional forman parte de esta determinación, desarrollada
en \ref{sec:action}.

El funcional de Barbero--Immirzi evalúa una cadena orientada de doce
sectores y una costura de retorno sobre incidencias de cardinales 90
y 120. Se obtiene
\(\gamma_{\HMT}=0.2740076726944592747\ldots\), con paridad impar
bajo inversión angular. Las reglas sectoriales CKM producen, por su
parte, una matriz racional de determinante \(-3857/6\), tres ángulos
de mezcla y la fase \(\delta_{\rm deg}=9A_{\rm deg}\).
Su realización compleja es unitaria y determina el invariante de
Jarlskog. Ambos resultados reciben las coordenadas angulares ya
generadas: el funcional de Barbero no selecciona los ángulos de mezcla,
ni estos seleccionan retrospectivamente las reglas sectoriales.

La composición longitudinal precede a la evaluación gravitatoria.
La inscripción de las incidencias y el archivo complementario producen
\(r_N=5759/23040\). La regla longitudinal R4 aplica este coeficiente
a \(\alpha_{\HMT}^{16}L_*\) y determina \(L\); la regla métrica
R5 determina el operador radial positivo del mismo carácter que utiliza
la energía. La sección~\ref{g26:sec:rigidez-radial-es} enuncia ambas
composiciones y demuestra la unicidad de su realización conjunta con
la acción de fase, el reloj conjugado y la identificación de radio
reducido R8. La referencia \(L_*\), declarada igual a un metro en la
evaluación SI, conserva su función dimensional.

La relación entre Boltzmann y gravitación se concreta entonces en
\[
 k_B\Theta_{{\rm clk},a}\log3=\frac{h_a}{T_{108}},\qquad
 G_a=\frac{c_{\rm int}^3R_{108}^2}{\hbar_a},\qquad
 \frac{G_a\hbar_a}{c_{\rm int}^3}=\ell_P^2.
\]
Aquí \(R_{108}=L\) y \(t_0=\pi_{\HMT}L/(54c_{\rm int})\).
La primera igualdad determina la energía por decisión trítica y un
transductor térmico único una vez fijada la sección positiva de
temperatura. La segunda es la especialización circular del coeficiente
común determinado por la prueba radial. La tercera muestra
que la reciprocidad de acción y gravitación conserva la unidad de
área. Estas condiciones térmicas y geométricas permiten relacionar
los resultados sin identificar la determinación de una sección
dimensional con la elección de su coordenada numérica.

La distinción requiere atender a la naturaleza de cada objeto. Una razón
adimensional expresa una relación entre magnitudes. Una constante con
dimensiones ocupa una recta de cantidades y adquiere una coordenada numérica
al elegir una unidad. Una condición de frontera fija una realización dentro
de una colección de soluciones. Las tres determinaciones pueden intervenir
en una misma fórmula y conservan funciones diferentes en su demostración.
El debate de Duff, Okun y Veneziano sobre las constantes dimensionales
ilustra la importancia de explicitar estas elecciones
\cite{duffokunveneziano}. La formulación actual del Sistema Internacional
precisa, a su vez, la función de las constantes definitorias en la
realización de las unidades \cite{bipm2018,bipmbrochure}.

\subsection{Soporte aritmético, orientación y dinámica generadora}

En Holografía Modular Triádica, HMT, el soporte aritmético, la orientación
ternaria y la dinámica con memoria preceden a las lecturas escalares.
APP, \emph{All Possible Paths}, utiliza el soporte discreto
\(X_9=(\mathbb Z/9\mathbb Z)^2\), con el grafo de Cayley aditivo
generado por \(\pm(1,0)\) y \(\pm(0,1)\). Sobre cada posición
de su coordenatización \(1,\ldots,9\) actúan las dos evaluaciones, suma y
producto, conservando sus residuos y cocientes. TRIT asigna una firma
\(\{-1,0,+1\}\) que distingue orientación y régimen local.
TPK, \emph{Topological Phase Kernel}, compone selección, transporte,
actualización de acarreo y memoria, y prolongación compatible del estado.
El retorno de fase conserva el avance de memoria. Las operaciones y la
generación regional de \ref{sec:nucleo}--\ref{sec:eta} construyen las
regiones numéricas, el registro dodecafásico y alfa antes de su uso en
los lectores de acción, incidencia e información.
La representación analítica correlacionada de alfa prolonga a toda
profundidad la coordenada producida por la clausura dodecafásica.
La realización arquimediana determina los valores de cilindros compatibles;
la realización de incidencia conserva otra información del mismo origen.
Esta dependencia común permite estudiar cada constante como valor y
como estructura, y situar después la medición como contraste.

Esta distinción se materializa ya en la constitución del registro.
La evaluación de visitas y trazas orientadas de
\ref{subsec:registro-evaluacion-incidencial} produce tres recuentos enteros
por ventana, conserva sus residuos y cocientes y determina las veinticinco
coordenadas de \eqref{eq:registro-composicion-incidencial}. El dominio
incluye la colección de incidencias, sus multiplicidades y su orientación.
La proposición~\ref{prop:registro-36-residuos} caracteriza exactamente la
información residual que determina esa imagen; las transformaciones
de \ref{subsec:k-pi-h} y \ref{subsec:k-hadamard} permiten después recuperar
el registro firmado. Quedan enlazadas así la evaluación del libro y la
reconstrucción de su salida, con una función demostrativa propia para
cada operación.

Esta estructura interviene efectivamente en las ecuaciones.
Los cardinales de incidencia intervienen en
un funcional de retorno; las secciones de acción intervienen en una
reciprocidad gravitatoria; las ocupaciones de borde intervienen en el área
y en el espectro modular; la inversión de una forma cuadrática interviene
en el intercambio de momento y enrollamiento. Cada conexión se expresa
mediante un mapa, una identidad y un dominio de validez.

\subsection{Acción, incidencia e información de frontera}

El argumento propio del artículo II se concentra en la relación entre
acción, área, entropía y gravedad. El parámetro de Barbero--Immirzi ocupa
una posición central porque enlaza el funcional de incidencia con una
normalización areal y con su realización en la acción gravitatoria.
La formulación de Holst proporciona un lenguaje de realización posterior
para este coeficiente \cite{holst1996}. En la construcción aquí estudiada,
el coeficiente se evalúa desde la incidencia orientada y los canales
angulares ya obtenidos.

La inclusión de una hexada marcada en una octada conserva el par de
incidencia que determina los cardinales 90 y 120. La vuelta dodecafásica
distingue doce sectores de una frontera orientada de retorno. El lector
lineal de esa cadena produce las multiplicidades del funcional. El
coeficiente de polo y la evaluación decimal se calculan a continuación.
La secuencia revela por qué el valor de Barbero--Immirzi pertenece a una
estructura de orientación y de área, además de ocupar una posición en
una tabla de constantes.

El área y el Hamiltoniano modular describen aspectos complementarios
del borde. La primera lectura agrega ocupaciones. La segunda pondera
su procedencia sectorial. Para \(\Delta_{\rm mod}\ne0\) y
normalización conocida, el Hamiltoniano permite obtener el observable
ponderado \(D\); junto al total \(N\), este reconstruye las dos
ocupaciones que originan el agregado. La inversión de la pareja
\((N,D)\) no requiere esa condición adicional. Esta recuperabilidad precisa
el contenido de la tesis genealógica: dos configuraciones pueden
compartir una lectura escalar y conservar diferencias accesibles a otra
lectura del mismo estado. El argumento se desarrolla primero en el
espacio finito correspondiente y después en la realización declarada.

La conservación de memoria precisa también cómo reutilizar una dinámica
cuando se eliminan variables auxiliares. La identidad
\eqref{mem:identidad-finita} retiene el estado terminal de cada truncación;
la proposición~\ref{mem:prop-serie} proporciona la prolongación y su cota
de cola. La proposición~\ref{mem:prop-schur} demuestra que la observación
completa se conserva al transformar conjuntamente el operador reducido,
la fuente y el lector. Esta regla identifica el contenido operativo que
debe transmitirse entre dos descripciones de un mismo transporte.

La entropía de entrelazamiento requiere además una bipartición y un
estado. La exposición especifica ambos mediante una purificación y su
estado reducido. La información adimensional del espectro se distingue
de la entropía dimensional, cuya normalización incorpora la constante
de Boltzmann. Esta precisión permite relacionar la energía de la
periodicidad trítica con la acción sin confundir un recuento de estados,
una distribución probabilística y una temperatura.

\subsection{Capacidad, memoria energética y realización termométrica}

La sección~\ref{sec:ii-ter27-antecedentes} desarrolla el mecanismo finito
que enlaza renovación de las dos hojas APP, capacidad y primera vacancia,
doble graduación del portador y energía de memoria. El lazo bilateral
determina un espectro decimal; su agrupamiento trítico conserva el resto.
La marca de capacidad y los sectores de incidencia producen después
una intensidad completa con su origen energético registrado. Las corrientes
conectadas fijan los saltos entre sectores, y el transporte centrado
conserva el lector entre realizaciones con escalas distintas.

La sección~\ref{sec:ii-termica-marcada} compone estos resultados con
acción, área, Barbero y horizonte. Añade una realización termométrica
que conserva la referencia espectral en un registro tensorial: el canal
entrópico lee la información marcada y el energético la esperanza
condicionada. La prueba obtiene el coeficiente de su cociente diferencial
y demuestra el principio variacional correspondiente. La referencia fija
y el par de funcionales marcado--condicionado son las condiciones
constitutivas de la realización. El Hamiltoniano cronológico de capacidad
conserva los grados $23,26,29$ y determina su estado normalizado; la
aplicación térmica de la intensidad se enlaza así con el reloj y el horizonte,
manteniendo sus escalas y sus transportes efectivos.

\subsection{Composición longitudinal y determinación del acoplamiento gravitatorio}

La gravitación forma parte del argumento central. La composición de
inscripción, longitud y métrica fija primero \(L\) y la respuesta radial.
El reloj conjugado satisface \(\omega_{108}L=c_{\rm int}\); su período
se distribuye en los 108 pasos del calendario dodecafásico. De ahí
proceden \(T_{108}=2\pi_{\HMT}L/c_{\rm int}\) y
\(t_0=T_{108}/108\). La condición de radio reducido R8 determina
\(G\) sobre el mismo carácter que ya expresa la energía. La
coincidencia circular conserva la forma equivalente
\(\vartheta_{108}^{\circ}=(54/\pi_{\HMT})^2\). En las secciones
de acción resultantes, el coeficiente gravitatorio se transforma
recíprocamente y el producto de acción y gravitación permanece
invariante. La unidad areal es \(L^2\) en ambas orientaciones.

La igualdad conservada posee una consecuencia geométrica:
el paso entre las secciones afecta a determinadas coordenadas mientras
preserva una cantidad geométrica común. El significado del valor se
obtiene del transporte que lo produce. La relación con el horizonte
introduce después la energía, la temperatura y la entropía areal bajo
normalizaciones explícitas. La confluencia de esas cantidades en una
celda de acción se demuestra por sustitución de las construcciones
anteriores.

La dependencia térmica se presenta con igual precisión. Una identidad
para el producto de la constante de Boltzmann por una temperatura
determina una escala energética. La determinación de cada coordenada
individual exige su correspondiente coordenatización. La normalización de cada realización fija las bases y el transporte
dimensional de esa coordenada.
La comparación física se realiza sobre esas salidas y sus unidades.

\subsection{Reciprocidad, propagación angular y continuidad excepcional}

Las coordenadas angulares intervienen en dos realizaciones diferenciadas
de la misma construcción. El funcional de incidencia orientada determina
el parámetro de Barbero--Immirzi; el lector de especie determina los
ángulos de mezcla CKM y su fase. La sección~\ref{sec:ii-propagacion-angular}
expone la matriz racional, su inversa y la realización unitaria compleja.
El determinante de la transformación y el invariante de Jarlskog
proporcionan controles complementarios: el primero examina la recuperación
de las coordenadas; el segundo examina la fase de las amplitudes. Los
valores numéricos se contrastan después de esas construcciones.

El teorema~\ref{thm:ii-witt-angular-transport} realiza las coordenadas
angulares en el módulo de incidencia de Witt mediante una isometría
explícita. El corolario~\ref{cor:ii-witt-composed-reader} conserva la
métrica del sistema sectorial y proporciona una inversa sobre su imagen.
Por ello, el paso a la incidencia excepcional transmite las coordenadas
producidas por las reglas angulares, junto con sus relaciones métricas.
La evaluación sectorial y el transporte isométrico se componen en ese
orden, conforme a \eqref{eq:ii-witt-composed-reader}.

La elipse de acción introduce una forma de determinante fijo y un
parámetro de anisotropía. El intercambio de sus semiejes invierte el
radio modular. Dos transportes pueden conservar la misma forma de
energía y actuar con signos diferentes sobre la forma simpléctica.
La orientación de la acción distingue entonces transportes simplécticos
y antisimplécticos con la misma energía escalar. Esta propiedad fundamenta
la atención al transporte y a la memoria en la continuación del modelo.

La relación acción--covarianza precisa el contenido cuántico de esta
geometría. La proposición~\ref{prop:ii-covarianza-gaussiana} construye,
en la representación canónica y con el estado gaussiano allí especificados,
una covarianza de determinante \(\hbar^2/4\) que alcanza la igualdad de
Robertson--Schrödinger. La proposición~\ref{prop:ii-area-covarianza-reciprocidad}
identifica el contorno de acción con la elipse de dos desviaciones
estándar: sus semiejes son \(2\Delta Q\) y \(2\Delta P\), y su área
es \(h=2\pi\hbar\). La reciprocidad intercambia las dispersiones y
conserva ese determinante. El paso a posición y momento mantiene explícita
la coordenatización dimensional, según~\eqref{eq:ii-covarianza-dimensiones}.
Así quedan relacionadas la acción previamente construida, su distribución
angular y las fluctuaciones de una realización operatoria definida.

La sección~\ref{delta22:ii:transporte} demuestra la conservación
exacta de la acción cuadrática bajo transporte entre elipses, deriva
su término de conexión hamiltoniana y establece la condición de
positividad del generador. La versión conformemente simpléctica
conserva la acción normalizada y explicita la forma transportada.


El intercambio momento--enrollamiento prolonga esta reciprocidad en
un espacio reticular con radio. Su demostración conserva el cociente
de acarreo, el dominio del operador y su transformación unitaria.
La función modular correspondiente se evalúa sobre el radio así
transportado. La incidencia excepcional se desarrolla hasta Moonshine
como parte de la segunda realización de la estructura común. Las
conexiones entre estas realizaciones se presentan por sus mapas
concretos, permitiendo distinguir qué cambia y qué se conserva.

\Needspace{7\baselineskip}
La comparación de los sectores quark y leptónico exige seguir por separado
los antecedentes que el desarrollo angular y los operadores de especie
transmiten a CKM y PMNS. Cada transporte conserva su marco propio; esa
distinción permite precisar dónde actúan la conjugación de amplitudes,
la orientación de rutas y los invariantes de mezcla. El manuscrito
mantiene las diferencias de construcción al pasar de las coordenadas
estructurales a las realizaciones respectivas.

\subsection{Identidades de área, entropía y escala gravitatoria}

La relación entre información y gravitación se expresa mediante tres
determinaciones articuladas. La unidad areal satisface
$\ell_P^2=G\hbar/c^3$. La incidencia orientada organiza su resolución
elemental mediante $\gamma_{\HMT}a_0$. El estado reducido y la bipartición
determinan la información accesible, cuya expresión dimensional utiliza
$k_B$. En la colección sectorial desarrollada en
\ref{sec:ii-ecuacion-informacion}, estas dependencias toman la forma
\begin{equation}
 \langle\widehat{\mathcal A}\rangle
 =\gamma_{\HMT}a_0\frac{G\hbar}{c^3}\langle N\rangle,
 \qquad
 S=k_B\bigl(\log Z+\Delta_{\rm mod}\langle D\rangle\bigr).
 \label{eq:ii-intro-area-informacion}
\end{equation}
Las cantidades $N$ y $D$ son las dos lecturas de las ocupaciones de
borde; $Z$ normaliza el estado. Su mapa conjunto recupera esas
ocupaciones mediante~\eqref{eq:recover}. Las pruebas contiguas muestran
también cómo sus fluctuaciones gobiernan las respuestas del área y de
la entropía.

El paso siguiente conserva sus datos geométricos explícitos. En la
realización de horizonte de la sección~\ref{sec:confluencia},
$S_H/k_B=A_H/(4\ell_P^2)$ y $T_HS_H=E_\Lambda$.
El área microscópica, la entropía del estado reducido y la entropía
de horizonte mantienen sus definiciones y sus aplicaciones de realización.
Esta organización permite leer en continuidad la incidencia, la acción,
la información y la gravitación, precisando en cada tránsito el objeto
que se conserva y el observable que cambia.

La sección escalar gravitatoria posee además una realización tensorial.
El descenso de incidencia y el proyector de rango cinco de
\ref{sec:ii-pentadica} determinan un portador pentacomponente;
su entrelazador lo realiza como tensores simétricos sin traza antes de
la reducción transversal de rango dos. La orientación de sus componentes
y la observación reducida permanecen así diferenciadas del acoplamiento
común. El área microscópica y la información sectorial, construidas en
\ref{sec:area} y \ref{sec:ii-ecuacion-informacion}, mantienen sus
definiciones al componer estos transportes. La continuidad excepcional
de \ref{exc:doble-lectura} y \ref{exc:leech} conserva la incidencia
hasta Moonshine; la dualidad de \ref{sec:ii-dualidad-t} conserva la
acción unitaria y los dominios del operador bajo el intercambio de
momento y enrollamiento.
% END THERMAL_R03_SYNC_01
